\newpage
\phantomsection
\label{prf:complement-antitone-inclusion}
\LRAProofFor{thm:complement-antitone-inclusion}

\begin{remark*}[Return]
\hyperref[thm:complement-antitone-inclusion]{Return to Theorem}
\end{remark*}

\begin{theorem*}[Complement is Inclusion-Antitone]
Let complements be taken relative to a fixed universe $U$. If
$A\subseteq B\subseteq U$, then $B^c\subseteq A^c$.
\end{theorem*}

\begin{proof}[Professional Standard Proof]
\LRAProofBodyStart
TODO: professional standard proof for complement-antitone-inclusion.
\end{proof}

\begin{proof}[Detailed Learning Proof]
\LRAProofBodyStart
TODO: detailed learning proof for complement-antitone-inclusion.
\end{proof}

\begin{remark*}[Proof structure]
\end{remark*}

\begin{dependencies}
\begin{itemize}
  \item \hyperref[def:subset]{Subset}
  \item \hyperref[def:complement]{Complement}
\end{itemize}
\end{dependencies}

\clearpage
